\documentclass[11pt]{article}

\usepackage[margin=1in]{geometry}
\usepackage{amsmath,amssymb,amsthm}
\usepackage{graphicx}
\usepackage[hidelinks]{hyperref}

\newtheorem{theorem}{Theorem}
\newtheorem{lemma}{Lemma}
\newtheorem{remark}{Remark}

\newcommand{\norm}[1]{\left\lVert #1\right\rVert}
\newcommand{\un}{\mathrm{un}}
\newcommand{\Kseq}{\mathcal K_{\mathrm{sun}}}

\title{KB-Spaces Characterize Monotone Stability of\\
Sequentially Un-Compact Operators from \texorpdfstring{$c_0$}{c0}}
\author{Automated Math Discovery Run \texttt{fa\_banach\_001}}
\date{August 9, 2026}

\begin{document}
\maketitle

\begin{center}
\fbox{\parbox{0.91\textwidth}{
\textbf{Status.} Candidate full solution, likely valid pending human review.
The novelty search described below found no later explicit answer, but novelty
confidence is moderate because the proof is short and non-arXiv literature was
not exhaustively citation-indexed.}}
\end{center}

\section*{Source and open question}

Kandi\'c, Marabeh, and Troitsky define a bounded operator into a Banach
lattice to be \emph{sequentially un-compact} when the image of every bounded
sequence has an un-convergent subsequence. In Proposition 9.6 of
\cite{KMT}, they prove that the following property forces $X$ to be a
KB-space:
\[
 \text{if }T_n\in L(c_0,X)\text{ are sequentially un-compact and }
 T_n\uparrow T,\text{ then }T\text{ is sequentially un-compact}.
\]
Immediately afterward they ask whether the converse is true. The exact
source location is page 30.

\begin{center}
\includegraphics[width=\textwidth]{figures/open_problem_crop.png}
\end{center}

Here $T_n\uparrow T$ is understood in the usual operator order: the sequence
is increasing and has supremum $T$; in particular, $T-T_1$ is positive.
A Banach lattice is a \emph{KB-space} when every increasing norm-bounded
sequence in its positive cone converges in norm. A net $(x_\alpha)$
un-converges to $x$ when
\[
 \norm{|x_\alpha-x|\wedge u}\longrightarrow0
 \qquad (u\in X_+).
\]

\section*{Proof intuition}

The monotone operator sequence itself barely needs to be analyzed. Its first
term already has the required sequential compactness, while the positive gap
$T-T_1$ has domain $c_0$. Positivity turns the vectors $Se_k$ into a positive
series whose partial sums are uniformly bounded: they are the images of the
unit vectors $e_1+\cdots+e_N$ of $c_0$. The KB property makes these partial
sums converge in norm. Their tails then dominate the image of every unit-ball
tail in $c_0$, uniformly, so $S$ is the operator-norm limit of its finite-rank
coordinate truncations. Thus the whole positive gap is compact, and adding it
to $T_1$ preserves sequential un-compactness.

\section*{The positive-gap lemma}

\begin{lemma}\label{lem:positive-c0}
If $X$ is a KB-space, then every positive bounded operator
$S:c_0\to X$ is compact.
\end{lemma}

\begin{proof}
Let $(e_k)$ be the standard unit-vector basis of $c_0$, let $P_N$ be the
projection onto the first $N$ coordinates, and put
\[
 x_k=Se_k\in X_+,
 \qquad
 s_N=\sum_{k=1}^N x_k=S(e_1+\cdots+e_N).
\]
The sequence $(s_N)$ is increasing and
\[
 \norm{s_N}\leq \norm{S}\norm{e_1+\cdots+e_N}_\infty=\norm{S}.
\]
Since $X$ is a KB-space, there is $s\in X_+$ such that
$s_N\to s$ in norm. Moreover $s_N\leq s$ for every $N$: for fixed $N$,
$s_m-s_N\in X_+$ when $m\geq N$, and the positive cone is norm closed.

Fix $N$ and $a\in B_{c_0}$. For $m>N$, positivity of $S$ gives
\begin{align*}
 \bigl|S P_m(I-P_N)a\bigr|
 &\leq S\bigl|P_m(I-P_N)a\bigr| \\
 &\leq \sum_{k=N+1}^{m} x_k
   =s_m-s_N
 \leq s-s_N.
\end{align*}
As $m\to\infty$, $P_m(I-P_N)a\to(I-P_N)a$ in $c_0$. Continuity of $S$ and
of the lattice modulus, together with norm closedness of $X_+$, therefore
yields
\[
 \bigl|S(I-P_N)a\bigr|\leq s-s_N.
\]
Consequently,
\[
 \norm{S(I-P_N)}\leq \norm{s-s_N}\longrightarrow0.
\]
Each $SP_N$ has finite rank, so $S$ is compact.
\end{proof}

\section*{Full characterization}

\begin{theorem}\label{thm:main}
For a Banach lattice $X$, the following are equivalent.
\begin{enumerate}
\item $X$ is a KB-space.
\item Whenever $(T_n)$ is a sequence of sequentially un-compact operators in
  $L(c_0,X)$ and $T_n\uparrow T$, the operator $T$ is sequentially
  un-compact.
\end{enumerate}
\end{theorem}

\begin{proof}
The implication (2)$\Rightarrow$(1) is Proposition 9.6 of \cite{KMT}.

Suppose (1), and let sequentially un-compact operators $T_n:c_0\to X$
satisfy $T_n\uparrow T$. Then
\[
 S:=T-T_1\geq0.
\]
Lemma~\ref{lem:positive-c0} shows that $S$ is compact. In particular, $S$ is
sequentially un-compact, since every norm-convergent sequence is
un-convergent.

The sequentially un-compact operators form a linear subspace
\cite[Proposition 9.2]{KMT}. For completeness, if $(z_j)$ is bounded, first
pass to a subsequence along which $T_1z_j$ un-converges, and then to a further
subsequence along which $Sz_j$ un-converges. Addition is continuous for
un-convergence, so $(T_1+S)z_j=Tz_j$ un-converges along the final subsequence.
Thus $T$ is sequentially un-compact, proving (2).
\end{proof}

\begin{remark}[A slightly stronger form]
The proof does not use the intermediate operators $T_2,T_3,\ldots$.
If $X$ is a KB-space, $R:c_0\to X$ is sequentially un-compact, and
$T-R\geq0$, then $T$ is sequentially un-compact. Indeed, the positive gap is
compact by Lemma~\ref{lem:positive-c0}.
\end{remark}

\section*{Verification notes}

The main possible pitfall is to write the non-$c_0$ vector
$(1,1,\ldots)$ and apply $S$ to it. The proof does not do this: $s$ is
constructed as a norm limit in $X$, and the tail estimate is first proved for
finite truncations. Sign-changing unit-ball vectors are covered by the
positivity inequality $|Sz|\leq S|z|$. Passing to the limit uses only norm
continuity and the closed positive cone. Finally, the order hypothesis is used
only in the valid implication $T_1\leq T\Rightarrow T-T_1\geq0$; no norm
convergence of $T_n$ is assumed.

An independent line-by-line audit is included in \texttt{verification.md}.
Its verdict is \emph{likely valid, candidate full solution pending human
review}.

\section*{Bounded novelty check and limitations}

On August 9, 2026, the search covered the run's registry, solution, attempt,
and proof-gap indexes; exact web searches for the source wording and close
variants involving ``sequentially un-compact'', ``order closed'', $c_0$, and
``KB-space''; and a local full-text scan of all 43 ingested arXiv source papers
that cite the exact title of \cite{KMT}, through the 2026 corpus. No later
paper stating or proving Theorem~\ref{thm:main} was found.

The search did locate later resolutions of two \emph{different} questions.
The paper arXiv:2304.04189 proves norm-closedness of un-compact operators,
while arXiv:2404.15641 answers the band-lifting question. This helps
separate those literature-closed targets from the converse addressed here.
The present novelty assessment is nevertheless bounded, not definitive:
non-arXiv literature was not exhaustively citation-indexed, and the positive
$c_0$ lemma is elementary enough that it may exist elsewhere in a different
form.

\section*{Human-review recommendation}

Review as a full solution. The decisive point is the uniform tail estimate in
Lemma~\ref{lem:positive-c0}; once it is accepted, the missing implication is
immediate. A reviewer should also check whether this elementary positive-
operator observation has already been explicitly connected to Proposition
9.6 in non-arXiv literature.

\begin{thebibliography}{9}
\bibitem{KMT}
M. Kandi\'c, M. A. A. Marabeh, and V. G. Troitsky,
\emph{Unbounded norm topology in Banach lattices},
J. Math. Anal. Appl. \textbf{451} (2017), 259--279;
arXiv:1608.05489; doi:10.1016/j.jmaa.2017.01.041.
\end{thebibliography}

\end{document}
